\subsection*{原收款只标识退款来源}

取一条新的独立历史：平台账户初始 10000，原购买已经成功收款 3000，X1 使余额成为 7000。Alice 退回其中一部分，希望返还 1000。她后来又退回另一部分，也应返还 1000。两个金额和原因相同，业务上却是两项合法退款。

如果把 X1 当成退款操作的唯一键，两项合法退款会被合并。如果每次 HTTP 重试都新建退款身份，同一项 1000 又可能被退两次。因此 X1 表示来源，每项退款另有稳定身份。沿用完整四分量 $S$，退款动作取 \code{refund.create.v1}，登记会话取 ER1，键取 RK1。首次受理把这份 $S_{\rm RF}$ 绑定到实际退款编号 RF1。

原输入为：
\begin{BookCode}
source_charge_id=C1
currency=CNY
amount_minor=1000
reason=customer_request
\end{BookCode}
RF1 同时保存已核验的来源 X1、原输入解释版本、平台退款范围与 \code{refund-RF1}、原平台退款输入，以及取得终态后保存的原结果。来源核对通过 C1 的成功结果取得 X1，并确认 C1 属于当前账户及退款授权范围。

\CFifteenFigure{ch15-09-refund-identity}{X1 是共同来源，两项合法退款分别取得稳定身份、实际 RF 与平台退款结果}{fig:ch15-refund-identity}

实际平台也常把部分退款作为独立对象。Stripe 创建退款接口要求指定原 Charge 或 PaymentIntent，允许多次部分退款，总额以尚未退回的原收款为上限，amount 使用最小货币单位的正整数~\cite{stripe-refunds-create}。下面在这个业务问题上建立预占与恢复协议，使已受理的未知退款也有额度承担。

\subsection*{两个合法请求怎样争用同一份额度}

让原收款为 3000，两项不同退款各申请 2000，分别使用 RK1、RK2。若两个事务都先读取成功退款总额为 0，都判断 $2000\le3000$，再各自插入未决 RF，系统就接受了合计 4000 的退款责任。唯一 RF 防止同一项退款重复受理，原收款的累计约束则需要另一个共同仲裁位置。

本例让所有退款受理事务锁住同一来源 C1，再读取成功与未决占用总额。RF1 先取得锁，预占 2000 并提交。RF2 随后取得锁，看到可用额度只有 1000，返回本次退款额度不可用，没有建立新的退款关联。按来源分别仲裁的协议允许不同来源各自并行，针对同一来源的首次预占则按共同锁排序。

\CFifteenFigure{ch15-10-refund-concurrent}{两项各 2000 的退款争用原 3000，后取得来源锁的事务读取已经提交的预占}{fig:ch15-refund-concurrent}

下面给出首次受理与已有身份恢复的完整次序。身份比较在当前额度判断以前进行，同一 RF 的重试可以找回它已经取得的预占：

\begin{BookCode}
BEGIN
锁定原退款登记会话，核对授权、固定截止与活动资格
按完整 S_RF 查找已有绑定
若已有：
    按原版本解释并比较全部原输入与来源
    异参：返回 conflict，保持原绑定
    同参：返回原 RF 及原终态，或原预占与未决进度
若尚无：
    锁定原 C1，核验所属账户及可信成功 X1
    核验当前入口准入、退款输入、币种、正金额与动作资格
    读取 C1 下成功退款总额 R_succ 和未决占用总额 R_pend
    若 R_succ + R_pend + 本次金额 > 原收款额：
        返回 refund_capacity_unavailable，不创建退款
    分配从未使用的实际 RF
    固定原输入、X1、平台范围与 refund-RF
    共同保存未决退款、金额预占和持久恢复责任
COMMIT
提交确认后返回实际保存的 RF 与结果
\end{BookCode}

\code{refund\_capacity\_unavailable} 在这里描述本次首次受理没有取得额度。这个分支把 $S_{\rm RF}$ 留在尚未受理状态。另一项未决退款最终被拒绝、释放额度后，原来尚未受理的输入可以再次参加仲裁。已经建立 RF 的稳定拒绝则保存该项已受理退款的足够终结证据，后续同参恢复返回这个原结论。

同一 $S_{\rm RF}$ 的并发来件也要经过唯一仲裁。若不同入口可以独立锁住来源，却把身份绑定写到彼此不共享的表中，原 $S_{\rm RF}$ 会分裂。所有能创建 RF 的入口共同核对完整身份，并让身份绑定、来源预占与恢复待办处于同一权威提交。

下面的教学 SQL 将已有 RF 的同参恢复放在只读快照分支，先核验该快照中的准入资格。首次预占才按固定次序取得 g7 的共用目录锁，为连续决定位置排序，因此同一分片内不同来源的新受理也会串行。概念协议的额度仲裁按来源组织，这份最小实现另外付出了整分片首次受理串行的代价。

\subsection*{来源、退款与额度怎样保存在同一权威域}

沿用第~\ref{section:14-3}~节的教学布局，退款来源表由受控结果入口从可信 $R_J$ 建立。它保存原 C、所属账户、X、币种与原收款额，客户端提交的 X 或可退金额不能直接填入来源表。\code{result\_position} 指向该来源的原成功决定。退款表保存完整身份、输入版本和原输入，原金额与平台身份一旦受理便固定：

% SQL15-REFUND-DDL
\begin{BookCode}
CREATE TABLE j.refund_sources(
 c text PRIMARY KEY REFERENCES j.results(c),
 g text NOT NULL, result_position bigint NOT NULL,
 account_id bigint NOT NULL, x text NOT NULL UNIQUE,
 currency text NOT NULL, amount bigint NOT NULL CHECK(amount>0),
 FOREIGN KEY(g,result_position) REFERENCES j.decisions(g,position));
CREATE TABLE j.refunds(
 rf text PRIMARY KEY, g text NOT NULL REFERENCES j.directory(g),
 s jsonb NOT NULL UNIQUE CHECK(j.valid_s(s) IS TRUE),
 source_c text NOT NULL REFERENCES j.refund_sources(c),
 source_x text NOT NULL, input_version text NOT NULL,
 input jsonb NOT NULL, amount bigint NOT NULL CHECK(amount>0),
 status text NOT NULL CHECK(status IN ('pending','succeeded','rejected')),
 provider_scope text NOT NULL, provider_key text NOT NULL,
 provider_input jsonb NOT NULL,
 provider_result jsonb, result jsonb,
 UNIQUE(provider_scope,provider_key));
\end{BookCode}

下面给出首次预占和已有身份恢复的 SQL 入口，隔离级别取 PostgreSQL 的 READ COMMITTED。已有 RF 的恢复用同一条查询的快照读取目录、会话与原退款，先核对代次、开放状态、活动会话与截止时间，再比较全部原输入。快照读到的会话已经关闭时，原 RF 也不会使客户端越过准入。查询以后发生的并发关闭，排在这次读取之后；这条只读恢复的线性化点就是该查询的快照。已有结果和未决进度都按该快照返回。

没有找到 RF 时，写入分支按目录、会话、来源的固定顺序取锁，再重新检查资格与已有身份。登记、关闭、冻结和其他退款写入者遵守相同顺序，成功/拒绝合流也先锁目录再锁来源。这样不会出现持有来源锁再反向等待目录的环。已有身份的读取不持有分片目录排他锁，首次预占为分配连续决定位置付出整分片串行的代价。会话判期使用可信数据库时钟，需要区间时间源的部署改用前节的可信上界。授权者通过受限权限调用函数，从认证上下文构造 S，函数继续检查四个分量的规范类型，并按会话重建相同表示。

% SQL15-RESERVE
\begin{BookCode}
CREATE FUNCTION j.reserve_refund(
 p_g text,p_epoch bigint,p_s jsonb,p_input jsonb)
 RETURNS j.refunds LANGUAGE plpgsql AS $$
DECLARE d j.directory%ROWTYPE; e j.sessions%ROWTYPE;
        src j.refund_sources%ROWTYPE; r j.refunds%ROWTYPE;
        used_amount bigint; x_amount bigint; pos bigint; r_seq bigint;
        canonical_s jsonb; gate record; n bigint;
BEGIN
 IF j.valid_s(p_s) IS DISTINCT FROM true THEN
   RAISE EXCEPTION 'invalid refund identity encoding';
 END IF;
 SELECT dir.route_epoch,dir.mode,ses.account_id,ses.action,
        ses.session_id,ses.state,ses.deadline,to_jsonb(original.*) AS refund
 INTO STRICT gate
 FROM j.directory dir JOIN j.sessions ses ON ses.session_id=p_s->>2
 LEFT JOIN j.refunds original ON original.s=p_s WHERE dir.g=p_g;
 IF gate.route_epoch<>p_epoch OR gate.mode<>'open' THEN
   RAISE EXCEPTION 'admission closed';
 END IF;
 IF gate.state<>'active' OR clock_timestamp()>=gate.deadline
    OR gate.account_id IS DISTINCT FROM (p_s->>0)::bigint
    OR gate.action<>'refund.create.v1'
    OR p_s->>1 IS DISTINCT FROM gate.action THEN
   RAISE EXCEPTION 'invalid refund session';
 END IF;
 canonical_s:=jsonb_build_array(gate.account_id,gate.action,
                               gate.session_id,p_s->>3);
 IF canonical_s IS DISTINCT FROM p_s THEN
   RAISE EXCEPTION 'noncanonical refund identity';
 END IF;
 IF gate.refund IS NOT NULL THEN
   r:=jsonb_populate_record(NULL::j.refunds,gate.refund);
   IF r.g<>p_g OR r.input_version<>'refund-v1'
      OR r.input IS DISTINCT FROM p_input THEN
     RAISE EXCEPTION 'refund input conflict';
   END IF;
   RETURN r;
 END IF;
\end{BookCode}

快照中尚无原退款时，进入固定锁序的写入分支；等待锁以后重新取得目录和会话，并再次查找原身份：

% SQL15-RESERVE-CONT1
\begin{BookCode}
 SELECT * INTO STRICT d FROM j.directory WHERE g=p_g FOR UPDATE;
 IF d.route_epoch<>p_epoch OR d.mode<>'open' THEN
   RAISE EXCEPTION 'admission closed';
 END IF;
 SELECT * INTO STRICT e FROM j.sessions
  WHERE session_id=p_s->>2 FOR UPDATE;
 IF e.state<>'active' OR clock_timestamp()>=e.deadline
    OR e.account_id IS DISTINCT FROM (p_s->>0)::bigint
    OR e.action<>'refund.create.v1' OR p_s->>1 IS DISTINCT FROM e.action THEN
   RAISE EXCEPTION 'invalid refund session';
 END IF;
 canonical_s:=jsonb_build_array(e.account_id,e.action,e.session_id,p_s->>3);
 IF canonical_s IS DISTINCT FROM p_s THEN
   RAISE EXCEPTION 'noncanonical refund identity';
 END IF;
 SELECT * INTO r FROM j.refunds WHERE s=canonical_s;
 IF FOUND THEN
   IF r.g<>p_g OR r.input_version<>'refund-v1'
      OR r.input IS DISTINCT FROM p_input THEN
     RAISE EXCEPTION 'refund input conflict';
   END IF;
   RETURN r;
 END IF;
\end{BookCode}

首次预占再锁来源，比较累计责任，分配实际 RF 并共同保存原受理决定：

% SQL15-RESERVE-CONT2
\begin{BookCode}
 SELECT * INTO STRICT src FROM j.refund_sources
  WHERE c=p_input->>'source_charge_id' FOR UPDATE;
 x_amount:=(p_input->>'amount_minor')::bigint;
 IF src.g<>p_g OR src.account_id<>e.account_id
    OR p_input->>'currency' IS DISTINCT FROM src.currency
    OR x_amount IS NULL OR x_amount<=0
    OR p_input->>'reason' IS NULL THEN
   RAISE EXCEPTION 'invalid refund input or source';
 END IF;
 SELECT COALESCE(sum(amount),0) INTO used_amount FROM j.refunds
  WHERE source_c=src.c AND status IN ('pending','succeeded');
 IF x_amount>src.amount OR used_amount>src.amount-x_amount THEN
   RAISE EXCEPTION 'refund_capacity_unavailable';
 END IF;
 UPDATE j.directory
  SET next_refund=next_refund+1,last_position=last_position+1
  WHERE g=p_g AND route_epoch=p_epoch AND mode='open'
  RETURNING next_refund,last_position INTO r_seq,pos;
 GET DIAGNOSTICS n=ROW_COUNT;
 IF n<>1 THEN RAISE EXCEPTION 'refund allocation count %',n; END IF;
 INSERT INTO j.refunds VALUES(
  'RF'||p_epoch||':'||r_seq,p_g,canonical_s,src.c,src.x,
  'refund-v1',p_input,x_amount,'pending',
  'simulator:merchant-7:test:refund.v1',
  'refund-RF'||p_epoch||':'||r_seq,
  jsonb_build_object('source_provider_charge_id',src.x,
    'currency',src.currency,'amount_minor',x_amount),NULL,NULL)
 RETURNING * INTO r;
 INSERT INTO j.decisions VALUES
  (p_g,pos,p_epoch,'refund_accept',to_jsonb(r));
 RETURN r;
END $$;
\end{BookCode}

调用者以一个事务执行函数，提交确认后返回实际保存的 RF。本例的新编号可能是 \code{RF8:1}，它对应前面协议中的第一项 RF1。异常使整个事务回滚，额度不足分支没有新 RF、没有新位置。同参恢复先于当前额度判断，取回原 RF 与原未决占用；输入或来源改变时完整比较报告冲突。\code{pending} 行本身就是持久恢复待办，扫描器按原范围和键继续，不以客户端连接存活作为条件。

结果合流先取可信成功或足够最终拒绝证据，再锁目录与来源，核对原 RF 的来源、金额和平台身份。下面展示已核验成功后的条件写入，演示限定第 8 路由代次的开放资格，冻结或换代即从失败出口回滚；实际入口把 \code{RF8:1} 替换为读取的原 RF，把完整原平台响应和接口原结果保存下来：

% SQL15-RESULT
\begin{BookCode}
BEGIN;
DO $$
DECLARE r j.refunds%ROWTYPE; d j.directory%ROWTYPE;
        src j.refund_sources%ROWTYPE; pos bigint; n bigint;
BEGIN
 SELECT * INTO STRICT d FROM j.directory WHERE g='g7' FOR UPDATE;
 IF d.route_epoch<>8 OR d.mode<>'open' THEN
   RAISE EXCEPTION 'stale refund result authority';
 END IF;
 SELECT * INTO STRICT src FROM j.refund_sources WHERE c='C1' FOR UPDATE;
 IF src.g<>d.g OR src.x<>'X1' OR src.currency<>'CNY' THEN
   RAISE EXCEPTION 'refund source ownership conflict';
 END IF;
 SELECT * INTO STRICT r FROM j.refunds WHERE rf='RF8:1' FOR UPDATE;
 IF r.g<>d.g OR r.status<>'pending' OR r.source_c<>'C1'
    OR r.source_x<>'X1' OR r.amount<>1000
    OR r.provider_scope<>'simulator:merchant-7:test:refund.v1'
    OR r.provider_key<>'refund-RF8:1' THEN
   RAISE EXCEPTION 'unexpected original refund';
 END IF;
 UPDATE j.refunds SET status='succeeded',
  provider_result='{"refund_id":"PX-RF1","source":"X1",
                    "currency":"CNY","amount_minor":1000}',
  result='{"refund_id":"RF8:1","source_charge_id":"C1",
           "source_provider_charge_id":"X1","status":"succeeded",
           "provider_refund_id":"PX-RF1","currency":"CNY",
           "amount_minor":1000}'
 WHERE rf=r.rf AND status='pending' AND input=r.input
 RETURNING * INTO r;
 GET DIAGNOSTICS n=ROW_COUNT;
 IF n<>1 THEN RAISE EXCEPTION 'refund result update count %',n; END IF;
 UPDATE j.directory SET last_position=last_position+1
  WHERE g=d.g AND route_epoch=d.route_epoch AND mode='open'
  RETURNING last_position INTO pos;
 GET DIAGNOSTICS n=ROW_COUNT;
 IF n<>1 THEN RAISE EXCEPTION 'refund decision count %',n; END IF;
 INSERT INTO j.decisions VALUES
  (d.g,pos,d.route_epoch,'refund_result',to_jsonb(r));
END $$;
COMMIT;
\end{BookCode}

已有终态的重复证据走另外的恢复分支：逐字段比较原两类结果，相同就返回原结果，不再写一项新决定。相反终态或字段冲突保留原决定并进入核查。最终拒绝只有在原退款的过去、全部已受理与在途、未来路径都取得相应终结依据后，才把同一行从 \code{pending} 改为 \code{rejected}。状态改变与完整结果决定共同提交，使扫描器停止、额度释放和后续重试使用同一事实。

\subsection*{本地退款投影按原决定恢复}

W 中另建只读退款投影，列与权威退款行对应。普通查询通过 RF 或完整 S 读取，首次预占、额度汇总和平台发送都由 J 中的原退款承担。只有受限恢复角色可以调用下面的物化入口，业务查询角色只有读取权限。投影缺失或回退时，恢复者按同一分片的决定顺序重放 \code{refund\_accept} 与 \code{refund\_result}，中间其他种类的决定也各自物化并推进同一游标。

% SQL15-W-REFUND-DDL
\begin{BookCode}
CREATE TABLE w.refunds(
 rf text PRIMARY KEY,g text NOT NULL,s jsonb NOT NULL UNIQUE,
 source_c text NOT NULL,source_x text NOT NULL,
 input_version text NOT NULL,input jsonb NOT NULL,
 amount bigint NOT NULL CHECK(amount>0),
 status text NOT NULL CHECK(status IN ('pending','succeeded','rejected')),
 provider_scope text NOT NULL,provider_key text NOT NULL,
 provider_input jsonb NOT NULL,provider_result jsonb,result jsonb,
 CHECK(j.valid_s(s) IS TRUE),
 UNIQUE(provider_scope,provider_key),
 CHECK((status='pending' AND provider_result IS NULL AND result IS NULL)
    OR (status IN ('succeeded','rejected')
        AND provider_result IS NOT NULL AND result IS NOT NULL)));
REVOKE ALL ON w.refunds FROM PUBLIC;
\end{BookCode}

物化参数 \code{p\_decision} 是已核验来源的完整决定，包含原分片 \code{g}、位置 \code{position}、代次 \code{route\_epoch}、类型 \code{kind} 和行图像 \code{body}。\code{p\_epoch} 则是本地当前恢复资格，两者的代次可以不同：第 9 代的恢复者会重放第 8 代的原退款。终态推进还携带同一 RF 的原受理行图像 \code{p\_previous}，只允许改变状态和两类结果。源 C 的成功、X、币种、金额及账户要先与本地已经恢复的原付款对应：

% SQL15-W-REFUND
\begin{BookCode}
CREATE FUNCTION w.materialize_refund(
 p_g text,p_epoch bigint,p_position bigint,p_decision jsonb,
 p_previous jsonb DEFAULT NULL) RETURNS bigint
 LANGUAGE plpgsql AS $$
DECLARE r w.refunds%ROWTYPE; old_r w.refunds%ROWTYPE;
        c w.facts%ROWTYPE; p w.facts%ROWTYPE; n bigint;
        image jsonb; old_image jsonb; decision_kind text;
BEGIN
 PERFORM 1 FROM w.cursor WHERE g=p_g AND accepted_epoch=p_epoch
   AND mode='recovering' AND last_position=p_position-1 FOR UPDATE;
 IF NOT FOUND THEN RAISE EXCEPTION 'invalid refund recovery gate'; END IF;
 IF p_decision->>'g' IS DISTINCT FROM p_g
    OR p_decision->>'position' IS DISTINCT FROM p_position::text
    OR p_decision->>'route_epoch' IS NULL THEN
   RAISE EXCEPTION 'refund decision source conflict';
 END IF;
 decision_kind:=p_decision->>'kind';
 image:=p_decision->'body';
 r:=jsonb_populate_record(NULL::w.refunds,image);
 IF image IS DISTINCT FROM to_jsonb(r) OR r.g IS DISTINCT FROM p_g
    OR (decision_kind='refund_accept' AND r.status<>'pending')
    OR (decision_kind='refund_result'
        AND r.status NOT IN ('succeeded','rejected'))
    OR decision_kind IS NULL
    OR decision_kind NOT IN ('refund_accept','refund_result') THEN
   RAISE EXCEPTION 'invalid refund decision image';
 END IF;
\end{BookCode}

完整原图像通过编码检查以后，物化者从本地原付款核对退款来源，再比较已存退款：

% SQL15-W-REFUND-CONT1
\begin{BookCode}
 SELECT * INTO STRICT c FROM w.facts
  WHERE kind='C' AND id=r.source_c;
 SELECT * INTO STRICT p FROM w.facts
  WHERE kind='P' AND id=c.parent_id;
 IF c.body->>'state' IS DISTINCT FROM 'succeeded'
    OR c.body->'provider_result'->>'provider_charge_id'
       IS DISTINCT FROM r.source_x
    OR c.body->'provider_result'->>'currency'
       IS DISTINCT FROM r.input->>'currency'
    OR c.body->'provider_result'->>'amount_minor' IS NULL
    OR (c.body->'provider_result'->>'amount_minor')::bigint<r.amount
    OR p.body->'s'->0 IS DISTINCT FROM r.s->0 THEN
   RAISE EXCEPTION 'refund projection source conflict';
 END IF;
 SELECT * INTO old_r FROM w.refunds WHERE rf=r.rf FOR UPDATE;
 IF NOT FOUND THEN
   IF decision_kind<>'refund_accept' THEN
     RAISE EXCEPTION 'missing original refund acceptance';
   END IF;
   INSERT INTO w.refunds SELECT r.*;
 ELSIF to_jsonb(old_r)=image THEN
   NULL;
 ELSE
   old_image:=to_jsonb(old_r);
   IF decision_kind<>'refund_result'
      OR old_image IS DISTINCT FROM p_previous
      OR old_r.status<>'pending'
      OR (old_image-ARRAY['status','provider_result','result'])
         IS DISTINCT FROM (image-ARRAY['status','provider_result','result']) THEN
     RAISE EXCEPTION 'refund projection full image conflict';
   END IF;
   UPDATE w.refunds AS existing SET status=r.status,
     provider_result=r.provider_result,result=r.result
    WHERE existing.rf=r.rf AND to_jsonb(existing.*)=p_previous;
   GET DIAGNOSTICS n=ROW_COUNT;
   IF n<>1 THEN RAISE EXCEPTION 'refund projection update count %',n; END IF;
 END IF;
 -- Empty purchase/ledger batches still use the same atomic cursor gate.
 RETURN w.materialize(p_g,p_epoch,p_position,'[]','[]');
END $$;
REVOKE ALL ON FUNCTION w.materialize_refund(
 text,bigint,bigint,jsonb,jsonb) FROM PUBLIC;
\end{BookCode}

例如 J 在位置 3 保存 \code{RF8:1} 的原受理，位置 4 保存其成功，W 的位置 2 已经包含原成功 C1。恢复者先传入位置 3 的完整决定，再传入位置 4 的完整决定与位置 3 的原行图像。两次都保留 \code{RF8:1/refund-RF8:1}，第二次保存原 \code{PX-RF1} 及全部原结果。W 回退到位置 2 以后重复这两步，得到相同的退款行，J 的预占、编号和决定位置没有变化。任何字段冲突、来源错配或游标零行，都回滚该次投影与游标。提交未知时先重读游标与完整原行，已越过的位置核对后确认，尚未越过的位置沿原决定继续。


\subsection*{未知不能释放已经接受的责任}

继续使用原收款 3000。RF1 预占 1000，平台按 \code{refund-RF1} 成功返还，余额从 7000 增为 8000，成功响应却丢失。J 中 RF1 未决，占用 1000。恢复者依据原平台范围与退款键取回同一项结果，核对 RF1、原 C1/X1、币种和金额，再把 RF1 从未决改为成功。

这次迁移让 $R_{\rm pend}$ 减 1000、$R_{\rm succ}$ 增 1000，$R_{\rm succ}+R_{\rm pend}$ 保持。若成功响应丢失便释放 $R_{\rm pend}$，另一项退款可能先把这 1000 用掉，原 RF1 的成功证据稍后到来，累计已退就会超过原收款。未知时保存占用，使本地尚未取得平台结果的期间也有额度承担。

对最终拒绝，恢复者核对原身份既往没有成功，已受理与在途路径都已终结，未来首次受理已经关闭。证据进入 RF1 的结果仲裁后，把 \code{pending} 改为 \code{rejected}，释放原预占。一次查询查空、暂时错误或工作人员取消本地任务时，原平台路径继续保留可能成功的责任。

\CFifteenFigure{ch15-11-refund-evidence}{未决退款保留占用，可信成功与足够最终拒绝分别决定两条终态出口}{fig:ch15-refund-evidence}

退款平台的所有恢复继续使用 \code{refund-RF1}。平台普通查询暂未看到退款时，恢复者继续核对原关联。此时改用 \code{refund-RF2}，会让两项退款都可能成功。若平台逐键保护到期，恢复者停止无法保证安全到达的发送，继续查询、对账或取得平台同域关闭依据。第~\ref{section:12-5}~节已完整推导同域关闭。退款平台若提供对应能力，就在原退款身份处关闭未来首次受理，核对过去与全部已受理、在途工作，并把未终结者的持久查询编号 H 写回原 RF 的恢复记录。原收款成功用于核验来源，退款受理、重复和终结则沿 RF 自己的记录恢复。

成功结果保存实际平台退款编号。例如 RF1 的原结果为：

\begin{BookCode}
{
  "refund_id": "RF1",
  "source_charge_id": "C1",
  "source_provider_charge_id": "X1",
  "status": "succeeded",
  "provider_refund_id": "PX-RF1",
  "currency": "CNY",
  "amount_minor": 1000
}
\end{BookCode}

再次退回另一部分 1000，客户端明确确认新意图，使用 RK2，取得 RF2 和另一个实际平台退款编号。两次合法退款的原结果分别保存，RF1 的同参重试继续返回 \code{PX-RF1}。两项成功合计 2000，尚可受理额度为 1000。

退款的金额关系和归纳放在第~\ref{section:15-3}~节，退款绑定与未决预占也作为 J 的权威事实跨本地回退保存。

实际资金等式在同一历史里同步变化。原付款后账户为 7000，两项成功退款各 1000 后为 9000。第三项 RF3 预占 1000、平台结果未知时，真实账户是否已经成为 10000 要由平台事实决定。本地准入在尚未退款、以及已经退款但响应丢失这两种情况下，都把这项 1000 算入 $R_{\rm pend}$。RF3 最终拒绝后，账户与原资金事实保持，入口释放 1000 给以后新受理的 RF4。

退款把新的入账事实连接到原 X1。购买已付款、部分金额已退、退款有未决可以同时成为查询事实。商品退回、库存恢复和履约取消若有独立业务动作，各自保存授权、身份和结果，再由查询组合相应进度。

\subsection*{版本升级先恢复原来的解释}

现在让一个部署同时处理两种动作：
\begin{BookCode}
本地动作：purchase.create.v1
外部动作：purchase.create.external.v2
\end{BookCode}
旧本地输入包含商品、数量、报价、币种和金额，新外部输入另有 \code{payment\_method}，受理后产生 C/T 与平台身份。两种动作的 $S$ 不同，数据与效果也有不同归属。

取一条独立历史：本地账户初始 10000，旧 $S_{L1}$ 用 Q7 成功 3000，保存 P1/O1/D1 与原余额 7000。旧 $S_{L2}$ 随后成功 3000，当前本地余额成为 4000。服务部署外部 v2 以后，原 $S_{L1}$ 重试，首先依据保存的 \code{local-v1} 解释版本比较原输入，返回原 O1/D1 与 7000。新 v2 的 \code{payment\_method} 必填规则用于对应新动作的首次输入。

若统一用最新解析器要求旧来件补一个 \code{payment\_method}，再把它当作新外部购买，旧本地成功会多出一次平台付款。另一种错误把旧结果的 \code{balance\_after\_minor} 填成当前 4000，效果数量保持，原结果却改变。已有绑定分支先取得持久保存的输入解释版本与原结果格式，再按原语义恢复。

\CFifteenFigure{ch15-13-versions}{旧本地动作恢复原输入与 7000，新外部动作取得独立 C 与平台身份}{fig:ch15-versions}

输入升级可以增加新动作，也可以为同一动作增加明确的输入版本。后者需要规定新旧报文怎样比较。若 v2 只是给缺省字段填一个值，迁移者先证明这个值与旧语义等价，再提供确定的旧输入解码。若 v2 改变了价格、扣款来源或业务含义，系统保存旧版本并恢复原绑定，新意图通过新版本参与首次判断。

未知的旧外部 C1 又有单独的恢复分支。新部署读取旧 $A_J$ 中的原 \code{provider\_scope}、\code{pay-C1} 和支付输入，沿原版本适配器查询或安全重发。新适配器按保存的版本与字段规则解释旧平台结果。遇到无法解释的新枚举或不相容字段时，保留原证据进入核查，付款发送身份保持原关联。

平台接口版本与本书动作版本也分别保存。平台可能根据请求头决定响应形状，回调端点也可能采用另一份固定版本。Stripe 文档给出 \code{Stripe-Version} 请求头与回调端点的版本设置~\cite{stripe-versioning}\footnote{请求头可以覆盖账户默认 API 版本，回调事件可在端点创建时固定版本。本书据此保存适配器需要解释的版本与原报文，跨版本输入的业务等价关系则由各动作的输入规则确定。}。适配器保存自己需要解释的版本与原报文，规范化后才进入 C/R 的证据合流。

\subsection*{发布新读取器以后再改变保存格式}

假设旧服务只认识 \code{payment\_result-v1}，新服务准备增加 \code{refund\_summary}。迁移先让读取器能识别旧结果，并把新增摘要从已经核验的退款事实计算出来。原 \code{payment\_result} 中的 P/C/O/X、金额和成功结论保持，摘要表达当前累计退款与未决占用。这样一项旧成功可以连同新的退款进度被查询。

若必须改变结果保存格式，先部署同时理解新旧格式的服务与消费者，写入继续使用旧格式。发布者核对所有可能接收恢复查询、读取结果或重放事件的实例版本，确认它们都具备兼容能力，或已从路由和消费组隔离不兼容实例。预定的回退版本也在此时验证能够读取将要保存的新格式。两项检查都通过以后，才启用新格式写入。这一检查形成发布屏障。例如 V1 只能读旧格式，V2 已经写下新结果，查询若还会到达 V1，原成功就会暂时无法恢复。V1 把未知格式当成缺记录并进入首次购买分支时，还会把一次升级变成历史遗忘。

新格式写入启用以后，部署沿已经验证的兼容版本回退。兼容服务中的旧格式读取分支继续服务旧记录和原事件的恢复。迁移者核对历史记录已完成转换、原结果逐字段一致，并确认保留期内的查询和事件回放都有读取路径以后，才删除旧格式支持。此时退出的是兼容代码中的旧格式分支，所有可接流进程已经在写新格式以前跨过了兼容检查。

\CFifteenFigure{ch15-14-reader-migration}{全部可接流读取者兼容后启用新写入，旧格式分支按迁移和回放义务退出}{fig:ch15-reader-migration}

消息升级同样保留原 EV1 的身份与原版本。消费者用对应版本解码，再按原 I/B 取得 F501。确有新的业务授权，例如补发第二件商品，才建立独立补发身份。将旧事件换一个 \code{event\_id} 重放给新消费者，会改变接收去重观察，业务屏障 B 继续核对它是否已经授权过同一初次履约。
\FloatBarrier
